% ══════════════════════════════════════════════════════════════
%  SECTION 5: COMPLIANCE
% ══════════════════════════════════════════════════════════════

\section{Referee Compliance and Norm Convergence}
\label{sec:compliance}

The preceding sections operate at the match level. But the directive
was addressed to referees, and whether a regulatory instruction
``worked'' is ultimately a question about individual behavioral
change. This section documents compliance at the referee level,
establishing that the norm shift was not driven by compositional
change---new referees replacing old---but by nearly universal
individual adjustment.


% ──────────────────────────────────────────────────────────────
\subsection{Individual-Level Compliance}
\label{subsec:compliance_rates}
% ──────────────────────────────────────────────────────────────

Among the 85~referees who officiated at least 20~matches in both the
pre- and post-directive periods, every single one increased their
average second-half stoppage time. The compliance rate is 100\%, with
a mean within-referee increase of $+1.46$~minutes. Relaxing the match
threshold modestly reduces the rate: at $\geq 10$ matches in each
period, 98\% comply (100 of 102); at any number of pre- and
post-matches, 95\% comply (114 of 120).

We are explicit about these thresholds because a reader who encounters
``near-universal compliance'' and then discovers a lower rate in the
unrestricted sample would be right to feel misled. The 95\% figure
from the any-pre-and-post sample of 120~referees is the honest
headline number. The 100\% figure among high-match referees reflects
selection into a subsample with statistical power to detect
individual-level change---it tells us something real about experienced
referees, but it is not the population rate.

At the match level, 70.2\% of post-directive treated-league matches
exceeded the pre-directive league-specific median stoppage time. This
figure varies by league: 81.9\% in the Premier League, 79.9\% in the
Bundesliga, 68.3\% in Ligue~1, 65.7\% in La~Liga, and 56.4\% in
Serie~A. The Italian figure---barely above chance---is consistent with
the muted treatment effect documented in the bilateral decomposition.

The discrepancy between referee-level (95--100\%) and match-level
(70\%) compliance reflects different measurement constructs. A referee
whose \textit{average} stoppage time increases still officiate
individual matches below the pre-directive median---because
match-specific circumstances (few stoppages, no injuries, no VAR
reviews) sometimes produce shorter added time even under the new norm.
The referee-level rate captures the shift in each official's central
tendency; the match-level rate captures the overlap between the two
distributions.


% ──────────────────────────────────────────────────────────────
\subsection{Convergence}
\label{subsec:convergence}
% ──────────────────────────────────────────────────────────────

The coefficient of variation across referees fell from 0.173 pre-directive
(mean~$= 5.06$, SD~$= 0.87$, $N = 208$) to 0.144 post-directive
(mean~$= 6.47$, SD~$= 0.93$, $N = 156$). This is convergence in the
relative sense: referees became more similar in proportional terms, even
as the raw standard deviation rose from 0.87 to 0.93~minutes. The
higher mean absorbed and exceeded the increase in absolute dispersion.
What had been a wide range of individual conventions about ``how much
is enough'' compressed into a narrower band around a higher center.

The heterogeneity forest reveals the convergence mechanism. Among
Big~Five leagues, La~Liga referees shifted the most ($+2.03$~minutes),
followed by the Premier League ($+1.88$), Ligue~1 ($+1.70$), the
Bundesliga ($+1.69$), and Serie~A ($+0.84$). That Spanish referees
adjusted the most is consistent with their relatively low pre-directive
baseline: leagues with more room to move, moved more.

The regression-to-the-new-mean pattern extends to individuals. Referees
in the lowest tercile of pre-directive stoppage time increased by
$+1.65$~minutes, compared to $+1.55$ for the middle tercile and
$+1.50$ for the highest. Those who had previously added the least
adjusted the most, consistent with convergence toward a new focal point.
Experience tells a similar story: less experienced referees adjusted
more ($+1.56$~minutes) than medium-experience ($+1.28$) or
high-experience ($+1.30$) referees, though the differences are modest.


% ──────────────────────────────────────────────────────────────
\subsection{Speed and Persistence}
\label{subsec:persistence}
% ──────────────────────────────────────────────────────────────

Adaptation was immediate. The first matchweek of the 2023--24 season
showed a compliance rate of 83.3\% and mean stoppage of 6.63~minutes.
There was no learning curve, no gradual ramp-up, no period of
experimentation. When the first five matchweeks are pooled, the
compliance rate was 87.0\%; after matchweek 30, it was 70.9\%---a
slight late-season decline that may reflect end-of-season match dynamics
or selection effects as lower-profile matches are assigned to less
experienced officials.

The compliance data extend through three full post-directive seasons.
The average within-referee increase was $+1.26$~minutes in 2023--24,
$+1.31$ in 2024--25, and $+1.38$ in 2025--26. Far from decaying, the
shift has slightly \textit{increased} over time. This rules out a
Hawthorne interpretation in which referees temporarily adjusted under
perceived scrutiny and then reverted. The new norm is self-sustaining.


% ──────────────────────────────────────────────────────────────
\subsection{Interpretation: Coordination, Not Enforcement}
\label{subsec:compliance_interp}
% ──────────────────────────────────────────────────────────────

The near-universality and permanence of the compliance response speak
to the mechanism behind the pre-directive equilibrium. If referees had
genuinely preferred adding 3--4~minutes and resisted the directive out
of private conviction, we would observe substantial holdouts, gradual
erosion, or compensating behavior on other margins. We observe none.

The pre-directive variance reflected a \textit{coordination failure}:
heterogeneous conventions about the appropriate amount of added time,
sustained by the absence of a clear focal point, rather than genuine
preference heterogeneity. The directive resolved this by creating a
focal point \citep{rothockenfels2002}. Once a critical mass of
referees shifted, the new convention became self-reinforcing: a referee
adding only 3~minutes when the norm was 6 would face professional and
social pressure from assessors, players, coaches, and media even
without formal sanctions. The mechanism is closer to clinical guideline
adoption, where new practice norms diffuse rapidly once an authoritative
body issues a recommendation \citep{cabana1999compliance}, than to
enforcement-driven regulatory compliance, where behavior changes only
under monitoring and punishment \citep{duflo2013truth}.

This distinction matters for policy. If compliance required sustained
enforcement, one would expect decay as monitoring attention faded. The
absence of decay---indeed, the slight increase over three
seasons---suggests that FIFA altered the equilibrium itself, not merely
the payoffs within the old one. The World Cup provided the coordinating
event that crystallized expectations that had been building since VAR's
introduction but lacked a focal point to coalesce into a new norm.
Three seasons have confirmed that the shift is permanent.

